\documentclass[11pt,oneside]{article}
\usepackage[T1]{fontenc}
%\usepackage{lmodern}
\usepackage{amsmath,amssymb,amsthm}
\usepackage{booktabs}
\usepackage[margin=1.1in]{geometry}
\usepackage[hidelinks]{hyperref}

\newtheorem{proposition}{Proposition}[section]
\newtheorem{conjecture}[proposition]{Conjecture}
\newtheorem{question}[proposition]{Question}
\theoremstyle{remark}
\newtheorem{remark}[proposition]{Remark}

\newcommand{\Rp}{\mathbb{R}}
\newcommand{\C}{\mathbb{C}}
\newcommand{\re}{\operatorname{Re}}
\newcommand{\im}{\operatorname{Im}}

\title{A Toeplitz positivity test on the angular distribution of zeta zeros,\\
with an independent reproduction of the Connes--Consani--Moscovici\\ spectral realization}
\author{Von Jensen}
\date{August 2026}

\begin{document}
\maketitle

\begin{abstract}
The Keiper disk coordinate $z_\rho = 1 - 1/\rho$ places the nontrivial zeros of
$\zeta$ on the unit circle exactly when the Riemann Hypothesis holds. We equip
the resulting angular measure with its trigonometric moments and use finite
Toeplitz positivity as a numerical detector of hypothetical off-critical zeros.
Measured on the first $300$ zeros with injected off-line quadruples
$\rho = \tfrac12 + \delta + i\gamma$, the detector breaks positive
semidefiniteness at matrix order $N^{*}$ satisfying
$N^{*}\log(1/r) \in [0.085,\,0.34]$ across all tested $(\gamma,\delta)$, where
$r = |z_\rho|$; this is five to ten times below the scalar sensitivity threshold
$n \gtrsim \gamma^{2}/\delta$ of Voros for the Keiper--Li coefficients, and far
below the exponential scale $1/\log(1/r)$. The minimal eigenvalue of the
deviation-only Toeplitz matrix obeys the empirical law
$\lambda_{\min} \approx -0.2\,N^{3}\log^{2}(1/r)$. An on-circle control
injection never triggers detection. We relate the mechanism to the
Carath\'eodory--Fej\'er rigidity of nearly singular positive Toeplitz forms, the
same structure used constructively in the spectral realization of
Connes, Consani and Moscovici, and we report a full independent reproduction of
their central computation at $(\lambda, N) = (3, 120)$: all twenty published
eigenvalue--zero discrepancies are reproduced to every stated figure, and the
two hypotheses of their main theorem, simplicity and evenness of the minimal
Weil eigenvector, are verified numerically at this parameter point, with
$\varepsilon_N = 2.954\times 10^{-38}$ and spectral gap ratio
$1.4\times 10^{-7}$.
\end{abstract}

\section{Introduction}

Write $\rho = \beta + i\gamma$ for the nontrivial zeros of the Riemann zeta
function and
\begin{equation}\label{eq:keiper}
z_\rho \;=\; 1 - \frac{1}{\rho}\,,
\end{equation}
the conformal coordinate introduced by Keiper \cite{Keiper}, which maps the
critical line $\re s = \tfrac12$ to the unit circle. The Riemann Hypothesis is
the statement $|z_\rho| = 1$ for every $\rho$, and the Keiper--Li coefficients
\begin{equation}\label{eq:li}
\lambda_n \;=\; \sum_\rho \Bigl[\,1 - \Bigl(1 - \frac1\rho\Bigr)^{\!n}\,\Bigr]
\end{equation}
are nonnegative for all $n \ge 1$ if and only if RH holds \cite{Li}. Bombieri
and Lagarias \cite{BL} showed that this criterion is an instance of a general
positivity principle for multisets and connected it to Weil's quadratic
functional. Voros \cite{VorosSharp,VorosBook,VorosDisc} determined the sharp
large-$n$ dichotomy: under RH,
$\lambda_n \sim \tfrac n2(\log n + \gamma_E - 1 - \log 2\pi)$, while a zero at
$\tfrac12 + t + iT$ off the line contributes an exponentially growing
oscillation that becomes visible against the smooth part only for
$n \gtrsim T^2/|t|$, an uncertainty-principle barrier for the scalar sequence.

This note treats the zeros' angular positions on the Keiper circle as a measure
and interrogates that measure with the classical trigonometric moment problem.
Three results are reported, all numerical and all obtained with two independent
methods or explicit controls.

First, finite Toeplitz positivity of the empirical moments detects injected
off-critical zeros at matrix orders well below the scalar barrier. For an
injected quadruple at height $\gamma$ and displacement $\delta$, detection
occurs at order $N^{*}$ with $N^{*}\log(1/r) \in [0.085, 0.34]$ in every tested
configuration (Table~\ref{tab:detect}), where
$\log(1/r) = \tfrac12\log\bigl(1 + 8\delta/((2\delta-1)^2 + 4\gamma^2)\bigr)$ is
the radial displacement in the disk coordinate. The scalar barrier corresponds
to $n\log(1/r) \gtrsim 1$; the matrix test operates an order of magnitude below
it because an $N \times N$ Toeplitz form correlates all moments up to order $N$
simultaneously.

Second, the minimal eigenvalue of the Toeplitz matrix built from the deviation
moments alone grows cubically: $\lambda_{\min} \approx -0.2\,N^3 \log^2(1/r)$
over the measured range (Section~\ref{sec:law}). The quadratic factor is the
small-$n$ expansion of the off-circle deviation
$2\cosh(n\log r) - 2$; the cubic amplification in $N$ is a property of the
Toeplitz quadratic form that we state as a conjecture and pose as a problem in
the perturbation theory of Toeplitz matrices.

Third, the detection mechanism is the contrapositive of a rigidity phenomenon
that has recently been used constructively. The Carath\'eodory--Fej\'er
structure theorem \cite{CF} forces the kernel polynomial of a singular positive
Toeplitz matrix to have all its roots on the unit circle; the extension of this
principle by Connes and van Suijlekom \cite{CvS} underlies the spectral
realization of low zeta zeros by Connes, Consani and Moscovici \cite{CCM},
where a nearly null vector of the truncated Weil quadratic form generates, with
striking accuracy, the actual zeros. Section~\ref{sec:ccm} reports a complete
independent reproduction of the central computation of \cite{CCM} at
$(\lambda, N) = (3, 120)$, including numerical verification of the two standing
hypotheses of their main theorem at this parameter point. The reproduction also
isolates, through a deliberate perturbation, which part of their construction
guarantees reality of the approximating spectrum and which part carries the
arithmetic accuracy.

The detector is not proposed as a method for verifying RH in any range;
direct methods locate zeros far more efficiently. Its interest is
resolution-theoretic: it quantifies how much more a positivity criterion sees
when it is read as a matrix condition rather than a scalar sequence, on the
coordinate where the criterion is native.

\section{The angular form of the Keiper--Li framework}\label{sec:angular}

All statements in this section are elementary restatements of
\eqref{eq:keiper} and \eqref{eq:li} and are recorded for use in the sequel.
With $m(x) = (x-1)/(x+1)$ one has the identity
\begin{equation}\label{eq:mobius}
z_\rho \;=\; m(2\rho - 1),
\end{equation}
so the Keiper coordinate is the M\"obius compression of the zero centered at
the critical line. For an on-line zero $\rho = \tfrac12 + i\gamma$, $\gamma>0$,
the image lies on the unit circle at angle
\begin{equation}\label{eq:angle}
\theta_\gamma \;=\; \pi - 2\arctan(2\gamma) \in (0, \pi),
\end{equation}
and the conjugate zero contributes the angle $-\theta_\gamma$.

\begin{samepage}
\begin{proposition}\label{prop:sos}
If all zeros with $|\gamma| \le \Gamma$ lie on the critical line, their
contribution to $\lambda_n$ is
\[
\lambda_n^{(\Gamma)} \;=\; \sum_{0<\gamma\le\Gamma} 4\sin^{2}\!\Bigl(\frac{n\,\theta_\gamma}{2}\Bigr).
\]
\end{proposition}
\begin{proof}
For a conjugate pair on the circle, $1 - z^n + 1 - \bar z^{\,n} =
2 - 2\cos n\theta = 4\sin^2(n\theta/2)$, and \eqref{eq:angle} gives the angle.
\end{proof}
\end{samepage}

\begin{samepage}
\begin{proposition}\label{prop:rapidity}
Let $w_\rho = \tfrac12 \operatorname{Log}(2\rho - 1)$ with the principal branch.
Then $\rho$ lies on the critical line if and only if
$\im w_\rho = \pm\pi/4$.
\end{proposition}
\begin{proof}
$\re(2\rho - 1) = 0$ if and only if $\operatorname{Log}(2\rho-1) \in
\log|2\rho-1| + i\{\pm\pi/2\}$.
\end{proof}
\end{samepage}

Under $\operatorname{artanh}$, which inverts $m$, the coordinate $w_\rho$ is
the rapidity of the velocity $z_\rho$; RH states that all zero rapidities lie
on the two horizontal lines $\im w = \pm \pi/4$. Both propositions were
confirmed numerically on the first $300$ zeros: the identity \eqref{eq:mobius}
holds exactly in exact arithmetic, $\bigl||z_\rho| - 1\bigr| \le 1.3\times
10^{-26}$ at working precision $25$ digits, the angle formula \eqref{eq:angle}
agrees with $\arg z_\rho$ to $4.4 \times 10^{-26}$, and
$\im w_\rho - \pi/4$ vanishes identically.

\section{Two computations of $\lambda_n$}\label{sec:lambda}

The experiments below require accurate $\lambda_n$ for moderate $n$ together
with the zero angles. Two independent computations were run and compared.

Method A extracts $\lambda_{n}$ as Taylor coefficients of
$\frac{d}{dz}\log\xi\bigl(1/(1-z)\bigr)$ by discretized Cauchy integrals on the
circle $|z| = 0.75$ with $256$ nodes at $40$-digit working precision, using
$\xi'/\xi(s) = 1/s + 1/(s-1) - \tfrac12\log\pi + \tfrac12\psi(s/2) +
\zeta'/\zeta(s)$. No zeros enter this computation. The radius was chosen after
observing, at radius $0.9$, an aliasing contamination of size
$\lambda_{n+256}\,r^{256} \approx 10^{-9}$; at $r = 0.75$ the extracted
$\lambda_1$ agrees with the closed form $1 + \gamma_E/2 - \tfrac12\log 4\pi$ to
all fifteen printed digits. At fixed radius the extraction shares the
well-known cancellation wall of this subject: rigorous large-$n$ computation
requires on the order of $n$ bits of precision \cite{Johansson}, and the
present method is used only for $n \le 48$.

Method B evaluates Proposition~\ref{prop:sos} over the first $300$ zeros and
corrects the tail by the integral of the pair contribution
$2\re\bigl[1 - (1 - 1/\rho)^n\bigr]$ against the smoothed density
$\frac{1}{2\pi}\log\frac{\gamma}{2\pi}$.

\begin{table}[ht]
\centering
\begin{tabular}{rlll}
\toprule
$n$ & Method A (zero free) & Method B (zeros + tail) & A $-$ B \\
\midrule
1  & 0.023095708966121 & 0.023099086453 & $-3.4\times10^{-6}$ \\
2  & 0.092345735228047 & 0.092359245164 & $-1.4\times10^{-5}$ \\
5  & 0.575542714461177 & 0.575627151058 & $-8.4\times10^{-5}$ \\
13 & 3.81724005784795  & 3.81781082592  & $-5.7\times10^{-4}$ \\
21 & 9.61796072314292  & 9.61945000892  & $-1.5\times10^{-3}$ \\
48 & 40.8652694603672  & 40.8730461068  & $-7.8\times10^{-3}$ \\
\bottomrule
\end{tabular}
\caption{Keiper--Li coefficients by the two methods. The difference carries a
single sign and grows smoothly with $n$; it is the deterministic bias of
replacing the discrete zero sum beyond height $\gamma_{300} \approx 541.8$ by
the smoothed density, and it bounds the truncation systematics relevant to
Section~\ref{sec:detector}.}
\label{tab:lambda}
\end{table}

\section{The Toeplitz detector}\label{sec:detector}

Let $\mu_K = \sum_{k \le K}(\delta_{\theta_k} + \delta_{-\theta_k})$ be the
angular measure of the first $K$ zero pairs and
\begin{equation}
c_n \;=\; \int e^{-in\theta}\, d\mu_K \;=\; 2\sum_{k\le K}\cos\bigl(n\theta_k\bigr)
\end{equation}
its trigonometric moments. By the Herglotz--Toeplitz theorem
\cite{GrenanderSzego}, a two-sided sequence is the moment sequence of a
positive measure on the circle if and only if every finite Toeplitz matrix
$T_N = (c_{i-j})_{i,j=0}^{N-1}$ is positive semidefinite. Points on the circle
therefore keep every $T_N$ positive semidefinite regardless of RH
considerations, while a zero off the line, entering with its symmetric
quadruple $\rho$, $\bar\rho$, $1-\rho$, $1-\bar\rho$ and hence with velocities
$z, \bar z, 1/z, 1/\bar z$ at radius $r \ne 1$, contributes
\begin{equation}\label{eq:dev}
d_n \;=\; 2\cos(n\theta)\bigl(r^{\,n} + r^{-n}\bigr),
\end{equation}
which is unbounded and must eventually break positivity.

The detector is: fix $K = 300$, inject one quadruple at
$\rho = \tfrac12 + \delta + i\gamma$, and record the least $N^{*}$ at which the
smallest eigenvalue of $T_{N}$ falls below $-10^{-6} c_0$.

\subsection{Controls}

Injecting a unit-circle atom at the same angle as the off-line test point, with
the same total weight, never triggers detection: over matrix sizes up to
$3000$ the smallest eigenvalue stays above $-1.9\times10^{-10}$, four orders of
magnitude inside the tolerance. Detection therefore responds to off-circle
displacement and not to the presence of an added atom. The baseline matrices
are themselves numerically near-singular at all tested sizes, with
$\lambda_{\min}$ of order $10^{-10}$ against $c_0 = 600$: the angles
\eqref{eq:angle} accumulate at $0$ like $1/\gamma$, and clustered atoms make
the trigonometric moment problem nearly degenerate. This near-degeneracy is not
an obstacle but the operative mechanism; see Section~\ref{sec:ccm}.

The threshold dependence is weak: for $(\gamma,\delta) = (10, 0.2)$, sweeping
the tolerance from $10^{-3}$ to $10^{-8}$ moves $N^{*}$ only from $208$ to
$82$, an exponent of $0.081$ per decade. Detection is governed by the geometry
of the perturbed form against the near-null cluster space, not by
threshold crossing of a growing scalar.

\subsection{Resolution}

\begin{table}[ht]
\centering
\begin{tabular}{rrlrr}
\toprule
$\gamma$ & $\delta$ & $\log(1/r)$ & $N^{*}$ & $N^{*}\log(1/r)$ \\
\midrule
10    & 0.40 & $3.98\times10^{-3}$ & 83  & 0.331 \\
10    & 0.30 & $2.99\times10^{-3}$ & 91  & 0.272 \\
10    & 0.20 & $1.99\times10^{-3}$ & 111 & 0.221 \\
10    & 0.10 & $9.97\times10^{-4}$ & 140 & 0.140 \\
10    & 0.05 & $4.99\times10^{-4}$ & 171 & 0.085 \\
5     & 0.20 & $7.91\times10^{-3}$ & 27  & 0.214 \\
20    & 0.20 & $5.00\times10^{-4}$ & 673 & 0.336 \\
14.13 & 0.25 & $1.25\times10^{-3}$ & 154 & 0.193 \\
\bottomrule
\end{tabular}
\caption{Detection order for injected off-line quadruples. In every case
$N^{*}\log(1/r) < 0.34$: detection occurs while the exponential factor
$r^{-n}$ has grown by less than $40$ percent. The scalar visibility threshold
of \cite{VorosSharp,VorosBook} is $n\log(1/r) \gtrsim 1$, equivalently
$n \gtrsim \gamma^2/\delta$; for the row $(\gamma,\delta) = (10, 0.2)$ that
threshold is $n \approx 500$ against the measured $N^{*} = 111$, and the same
factor of five to ten separates the two across the table.}
\label{tab:detect}
\end{table}

The comparison in Table~\ref{tab:detect} is against a heuristic threshold
supported by saddle-point analysis rather than a theorem with constants, and it
should be read at that strength. The mechanism of the gain is nevertheless
transparent: a single coefficient $\lambda_n$ cannot resolve a radial
displacement until $n\,|{\log r}| \gtrsim 1$, which is the uncertainty
principle in the conjugate pair $(n, \log|z|)$ made explicit in
\cite{VorosBook}; the quadratic form on all moments up to $N$ is subject to no
such bound and instead feels the deviation \eqref{eq:dev} at its quadratic
onset $d_n - 4\cos(n\theta) \approx 2\cos(n\theta)\,n^2\log^2 r$.

\subsection{The eigenvalue law}\label{sec:law}

Isolating the deviation, the Toeplitz matrices built from
$d_n - 4\cos(n\theta)$ alone have smallest eigenvalues, at
$(\gamma, \delta) = (10, 0.2)$,
\[
-0.119,\; -0.790,\; -6.10,\; -46.3,\; -401
\quad\text{at}\quad N = 50, 100, 200, 400, 800,
\]
a factor close to $8$ per doubling.

\begin{samepage}
\begin{conjecture}\label{conj:law}
For the deviation symbol of a single off-circle quadruple at radius $r$ and
angle $\theta$, as $N \to \infty$ with $N\log(1/r) \to 0$,
\[
\lambda_{\min}\bigl(T_N[d - 4\cos(\cdot\,\theta)]\bigr)
\;=\; -\,c(\theta)\, N^{3} \log^{2}(1/r)\,\bigl(1 + o(1)\bigr),
\]
with $c(\theta)$ of order $0.2$ at the measured configuration.
\end{conjecture}
\end{samepage}

\begin{question}
Derive Conjecture~\ref{conj:law}, including the constant, from the asymptotic
spectral theory of Toeplitz matrices with polynomially growing symbol
coefficients \cite{BS}.
\end{question}

\section{Relation to the Carath\'eodory--Fej\'er mechanism, and a reproduction
of the spectral realization of Connes, Consani and Moscovici}\label{sec:ccm}

The structure theorem of Carath\'eodory and Fej\'er \cite{CF} states that a
positive semidefinite Toeplitz matrix with nontrivial kernel has a kernel
polynomial whose roots all lie on the unit circle. Read contrapositively, a
nearly singular positive Toeplitz form has no room for off-circle mass: any
such mass must break positivity against the near-null space essentially at
once, rather than after exponential amplification. This is the phenomenon
measured in Section~\ref{sec:detector}, where the clustered zero angles supply
the near-degeneracy and the injected quadruple supplies the off-circle mass.

The same structure is used in the opposite direction by Connes, Consani and
Moscovici \cite{CCM}, building on \cite{Connes99, CC23} and on the extension of
the Carath\'eodory--Fej\'er principle in \cite{CvS}. There, the input is
arithmetic rather than spectral: the Weil quadratic form restricted to test
functions on $[\lambda^{-1}, \lambda]$, involving only the prime powers up to
$\lambda^2$, is truncated to $2N{+}1$ Fourier modes; the eigenvector $\xi$ of
its smallest eigenvalue $\varepsilon_N$, made an exact kernel vector by
subtracting $\varepsilon_N$, turns a rank-one perturbation of the scaling
operator into an operator that is self-adjoint for the inner product defined by
the form. Its spectrum, which is the real zero set of the Fourier transform
$\widehat\xi$, is real by construction, and it approximates the actual ordinates
of the low zeta zeros to remarkable accuracy.

We reproduced the computation of \cite{CCM} at $(\lambda, N) = (3, 120)$
independently and from scratch: all closed-form matrix-element formulas were
re-derived and verified against direct quadrature before use, the matrix was
assembled through its structure
$\tau_{ij} = (b_i - b_j)/(i-j)$ for $i \ne j$, the problem was reduced by
parity, and the near-kernel eigenpair was extracted by inverse iteration with a
$200$-digit LU factorization. The constant
$c(L) = \int_0^L (e^{x/2}-1)/(e^x - e^{-x})\,dx = 0.575219384657911$ at
$L = 2\log 3$ enters every diagonal entry through $2c(L)$.

\begin{table}[ht]
\centering
\begin{tabular}{rll@{\qquad}rll}
\toprule
$k$ & this work & \cite{CCM} & $k$ & this work & \cite{CCM} \\
\midrule
1 & $1.58\times10^{-34}$ & $1.6\times10^{-34}$ & 11 & $7.28\times10^{-17}$ & $7.3\times10^{-17}$ \\
2 & $2.06\times10^{-31}$ & $2.1\times10^{-31}$ & 12 & $2.21\times10^{-15}$ & $2.2\times10^{-15}$ \\
3 & $1.46\times10^{-29}$ & $1.5\times10^{-29}$ & 13 & $9.75\times10^{-14}$ & $9.7\times10^{-14}$ \\
4 & $8.29\times10^{-27}$ & $8.3\times10^{-27}$ & 14 & $2.70\times10^{-13}$ & $2.7\times10^{-13}$ \\
5 & $1.28\times10^{-25}$ & $1.3\times10^{-25}$ & 15 & $6.34\times10^{-12}$ & $6.3\times10^{-12}$ \\
6 & $1.17\times10^{-23}$ & $1.2\times10^{-23}$ & 16 & $5.63\times10^{-11}$ & $5.6\times10^{-11}$ \\
7 & $7.53\times10^{-22}$ & $7.5\times10^{-22}$ & 17 & $2.93\times10^{-10}$ & $2.9\times10^{-10}$ \\
8 & $6.57\times10^{-21}$ & $6.6\times10^{-21}$ & 18 & $1.22\times10^{-9}$  & $1.2\times10^{-9}$ \\
9 & $1.19\times10^{-18}$ & $1.2\times10^{-18}$ & 19 & $5.58\times10^{-8}$  & $5.6\times10^{-8}$ \\
10 & $8.79\times10^{-18}$ & $8.8\times10^{-18}$ & 20 & $2.42\times10^{-7}$  & $2.4\times10^{-7}$ \\
\bottomrule
\end{tabular}
\caption{Absolute differences between the ordinates $\gamma_k$ of the first
twenty zeta zeros and the spectrum of the perturbed scaling operator at
$(\lambda, N) = (3, 120)$: the present reproduction against the values reported
in Figure~1 of \cite{CCM}. Agreement holds at every stated figure.}
\label{tab:ccm}
\end{table}

Two hypotheses of the main theorem of \cite{CCM} are assumed there for general
parameters: that $\varepsilon_N$ is simple and that $\xi$ is even. Both are
facts at this parameter point:
\[
\varepsilon_N \;=\; 2.954058093\times 10^{-38},
\qquad
\frac{\varepsilon_N}{\lambda_2^{\mathrm{even}}} \;=\; 1.38\times10^{-7},
\qquad
\lambda_{\min}^{\mathrm{odd}} \;=\; 1.13\times10^{-34},
\]
so the smallest eigenvalue is simple with seven orders of separation inside the
even sector and sits four orders below the entire odd sector. The size of
$\varepsilon_N$ is consistent with the superexponential decay displayed in
Figure~4 of \cite{CCM}.

\begin{samepage}
\begin{remark}\label{rem:split}
The reproduction separates the two guarantees of the construction. A
bookkeeping error in the eigenvector coordinates, a factor $\sqrt2$ between the
parity-reduced and the original basis, leaves every approximating eigenvalue
exactly real while destroying all arithmetic accuracy, the discrepancies in
Table~\ref{tab:ccm} inflating from $10^{-34}$ to order one. Reality of the
spectrum is structural and survives arbitrary corruption of $\xi$; proximity to
the Riemann zeros lives entirely in the fine structure of the near-kernel
vector. The detector of Section~\ref{sec:detector} and the realization of
\cite{CCM} are in this sense two readings of one rigidity: the near-null space
of a positive form on a finite window is, at once, the most sensitive locus for
off-line mass and the most information-dense encoding of the on-line zeros.
\end{remark}
\end{samepage}

\section{Scope and next steps}

The truncation systematics of the detector are bounded by
Table~\ref{tab:lambda}; the injected-quadruple experiments are synthetic, and
the planned validation is the Davenport--Heilbronn function, which satisfies a
functional equation while violating the analogue of RH and therefore supplies
genuine off-line zeros, as used for the discretized Keiper--Li sequences in
\cite{VorosDisc}. On the theoretical side the open items are
Conjecture~\ref{conj:law} and the extension of the detection-order analysis to
a statement with constants, for which the natural benchmark is the Oesterl\'e
bound quoted in \cite{VorosBook}, that verification of RH up to height $T_0$
forces $\lambda_n \ge 0$ for $n \le T_0^2$. The nearest published relatives of
the detector are the Fourier-positivity tests of Giraud and Peschanski
\cite{GP}, which watch the lowest eigenvalue of a Toeplitz hierarchy for a
candidate positive-definite function, and the model-space norm identity of
Suzuki \cite{Suzuki}, which realizes $\lambda_n$ as a squared norm; the present
test differs from the first in its object, the empirical angular measure of the
zeros, and from the second in being a finite-order matrix condition rather than
an identity.

\begin{thebibliography}{20}

\bibitem{Keiper} J.~B.~Keiper,
\emph{Power series expansions of Riemann's $\xi$ function},
Math. Comp. \textbf{58} (1992), 765--773.

\bibitem{Li} X.-J.~Li,
\emph{The positivity of a sequence of numbers and the Riemann hypothesis},
J. Number Theory \textbf{65} (1997), 325--333.

\bibitem{BL} E.~Bombieri and J.~C.~Lagarias,
\emph{Complements to Li's criterion for the Riemann hypothesis},
J. Number Theory \textbf{77} (1999), 274--287.

\bibitem{VorosSharp} A.~Voros,
\emph{A sharpening of Li's criterion for the Riemann hypothesis},
Math. Phys. Anal. Geom. \textbf{9} (2006), 53--63.

\bibitem{VorosBook} A.~Voros,
\emph{Zeta Functions over Zeros of Zeta Functions},
Lecture Notes of the Unione Matematica Italiana \textbf{8}, Springer, 2010.

\bibitem{VorosDisc} A.~Voros,
\emph{Discretized Keiper/Li approach to the Riemann hypothesis},
Exp. Math. \textbf{29} (2020), 452--469.

\bibitem{Coffey} M.~W.~Coffey,
\emph{Toward verification of the Riemann hypothesis: application of the Li
criterion}, Math. Phys. Anal. Geom. \textbf{8} (2005), 211--255.

\bibitem{AriasdeReyna} J.~Arias de Reyna,
\emph{Asymptotics of Keiper--Li coefficients},
Funct. Approx. Comment. Math. \textbf{45} (2011), 7--21.

\bibitem{Maslanka} K.~Ma\'slanka,
\emph{Li's criterion for the Riemann hypothesis: numerical approach},
Opuscula Math. \textbf{24} (2004), 103--114.

\bibitem{Johansson} F.~Johansson,
\emph{Rigorous high-precision computation of the Hurwitz zeta function and its
derivatives}, Numer. Algorithms \textbf{69} (2015), 253--270.

\bibitem{GrenanderSzego} U.~Grenander and G.~Szeg\H{o},
\emph{Toeplitz Forms and Their Applications},
University of California Press, 1958.

\bibitem{CF} C.~Carath\'eodory and L.~Fej\'er,
\emph{\"Uber den Zusammenhang der Extreme von harmonischen Funktionen mit
ihren Koeffizienten und \"uber den Picard--Landauschen Satz},
Rend. Circ. Mat. Palermo \textbf{32} (1911), 218--239.

\bibitem{BS} A.~B\"ottcher and B.~Silbermann,
\emph{Introduction to Large Truncated Toeplitz Matrices},
Springer, 1999.

\bibitem{Connes99} A.~Connes,
\emph{Trace formula in noncommutative geometry and the zeros of the Riemann
zeta function}, Selecta Math. (N.S.) \textbf{5} (1999), 29--106.

\bibitem{CC23} A.~Connes and C.~Consani,
\emph{Spectral triples and $\zeta$-cycles},
Enseign. Math. \textbf{69} (2023), 93--148.

\bibitem{CvS} A.~Connes and W.~D.~van Suijlekom,
\emph{Quadratic forms, real zeros and echoes of the spectral action},
Comm. Math. Phys. \textbf{406} (2025), 312.

\bibitem{CCM} A.~Connes, C.~Consani and H.~Moscovici,
\emph{Zeta spectral triples}, arXiv:2511.22755 (2025).

\bibitem{GP} B.~Giraud and R.~Peschanski,
\emph{From ``Dirac combs'' to Fourier-positivity},
Acta Phys. Polon. B \textbf{47} (2016), 1075--1100.

\bibitem{Suzuki} M.~Suzuki,
\emph{Li coefficients as norms of functions in a model space},
arXiv:2301.05779 (2023).

\end{thebibliography}

\end{document}
